\ifthenelse{\equal{\lang}{ko}}{%
\begin{abstract}
생성형 AI의 급격한 성장은 거버넌스와 컴플라이언스에 대한 수요를 강화하였다. 
기존의 모니터링 도구—수작업 보고서, 키워드 알림, 정적 추적기—는 
전 세계적으로 빠르게 변화하는 규제를 따라잡지 못한다. 
본 연구에서는 (1) 실시간 변화 탐지, (2) 의미론적 데이터 수집, 
(3) LLM 기반 보고를 통합한 \textbf{LexSense}라는 통합 프레임워크를 제안한다. 
새로운 분류 체계를 통해 다양한 관할권과 언어에서 규제 문서, 계약서, 소송, 
AI 자산 공개를 정밀하게 분류할 수 있다. 
EU, 미국, 한국 데이터셋을 대상으로 한 실험에서 91\%의 정확도, 
분석가 업무량 82\% 감소, 보고 속도 70\% 향상이라는 결과를 달성하였다. 
다국어 테스트와 소거(ablation) 실험을 통해 강건성이 확인되었다. 
또한 전체 코드, 데이터셋, Docker 기반 파이프라인을 공개하여 재현성을 보장한다. 
기술적 성과를 넘어, LexSense는 공정성 인식 튜닝, 편향 감사, 
프라이버시 보호형 모니터링을 통합하였다. 
이러한 기여는 \textbf{AI 거버넌스 정보학(AI Governance Informatics)}이라는 
새로운 연구 방향을 정립하며, 확장 가능하고 투명하며 책임 있는 
AI 거버넌스를 위한 배포 가능한 컴플라이언스 인프라와 개념적 토대를 함께 제공한다.
\end{abstract}



%\textbf{Keywords:}: 생성형 인공지능, AI 거버넌스 및 컴플라이언스, 대규모 언어 모델(LLM), 자동화 규제 모니터링, 책임 있는 AI 시스템


}{%
% -------------------- English Version --------------------
\abstract{The rapid growth of generative AI has intensified demands for governance and compliance. Existing monitoring tools---manual reports, keyword alerts, and static trackers---cannot keep pace with dynamic, global regulatory changes. We introduce \textbf{LexSense}, a unified framework combining (1) real-time change detection, (2) semantic data acquisition, and (3) LLM-driven reporting. A novel taxonomy enables fine-grained classification of governance documents, contracts, lawsuits, and AI asset releases across jurisdictions and languages. Experiments on EU, U.S., and Korean datasets achieve 91\% accuracy, 82\% less analyst effort, and 70\% faster reporting than manual baselines. Cross-lingual tests and ablations confirm robustness. We release full code, datasets, and Dockerized pipelines for reproducibility. Beyond technical gains, LexSense integrates fairness-aware tuning, bias audits, and privacy-preserving monitoring. Together, these contributions establish \textbf{AI Governance Informatics} as a new research direction, offering both a deployable compliance infrastructure and a conceptual foundation for scalable, transparent, and responsible AI governance. 
}


\keywords{AI Governance Informatics, Large Language Models (LLMs), Regulatory Monitoring, Compliance Automation, Cross-lingual Information Systems}
}
